\section{Unchecked inputs and remaining gates}
\label{sec:remaining}

The following entries are deliberately outside the checked-results
part. They are not assigned a proof status by analogy with a nearby
proved result.

\paragraph{U01: the requested Hooley-based proof.}
The strongest primitive-factor reduction is now A46: in every positive
characteristic and genus at least two it suffices to construct a pointed
odd cyclic etale cover, of degree prime to the characteristic, whose
actual primitive factor has no rational characteristic-primary torsion
over the specified finite field. A46 removes the previous elliptic
qualification and does not require geometric rank zero or a constant
deck group. Universal existence of such a cover is still unproved.

A02 proves the arithmetic consequence of M. A03 proves the exact
prime-degree geometric implication, and A11 weakens its sufficient
condition to an odd prime power $n$ with
$\langle2,-1\rangle=(\mathbb Z/n)^\times$ and an exact-order-$n$
theta point. A15 and A18 weaken this further: in characteristic two and every
genus at least two, a cyclic cover of any odd degree with primitive
factor of geometric $2$-rank zero suffices. In genus two this is
equivalent to theta meeting every primitive doubling orbit of the
defining cyclic torsion subgroup. A18 removes the exceptional
elliptic double quotient by an exact Deuring--Shafarevich and
conductor-rank calculation.
A04 refutes the inference from
two independent density statements, and A06 refutes the universal
joint primitive-root theta assertion itself. A05 proves A1 by a
different route. No claim is made here that a proof of A1 through
Hooley has been found. A strategy restricted to positive-two-rank
genus-two curves still requires a curve-dependent incidence theorem.
A07 refutes the positive-rank version over finitely generated fields
with transcendental moduli, not over finite fields. A08 shows that
no fixed finite prime list works over all finite fields; in the
ordinary case it excludes any fixed finite list of odd degrees.
A09 supplies an exact degree-$13$ rank-one witness. A13 supplies an
ordinary example of prime degree $4159$, with $2$ of index two in
the unit group. A16 supplies a rank-one example of degree $15$
whose full new factor has positive $2$-rank but whose primitive
factor has rank zero. A10 excludes order-three theta incidence on every nonordinary
genus-two base. None of
these results interchanges the finite-list and all-primes quantifiers;
in higher genus it requires nilpotence of every relevant Frobenius
orbit composite, not merely one nonordinary arrow.
A18 now permits any odd conductor with geometric primitive rank zero.
A19 makes all exponents for each specified finite odd prime support
an exact finite search. A21--A22 show that neither generic
positive rank nor any fixed finite support forces success, even
when mixed orders are allowed. A23 refutes universal success of
the older large-prime quotient inequality over constant extensions.
The remaining existence assertion must therefore allow the prime
support to depend on the particular finite-field curve, or use a
separating construction outside these sufficient tests.

A26 weakens the separate rational quotient criterion using even
augmentation. A27 nevertheless gives a rank-one curve over $\mathbf F_4$
where its numerical hypothesis fails for every extension and every
primitive-two prime. A28 recalls the recovered canonical-class criterion,
and A29 strictly improves its prime-level detection: a nonzero
$r$-primary component of $[K_C-2P]$ suffices if
$\operatorname{ord}_r(p)$ is even, with a cyclic $r^a$ cover and a
prime-to-$p$ rational top Prym. This is unconditional and needs no
prime-distribution theorem. The actual-curve strictness example in A29
shows why passage to prime powers matters even when all prime-level
canonical detectors fail. The residual points have canonical class
supported on the characteristic prime and odd primes of odd
multiplicative order, with the prime-two issue separate in odd
characteristic. In characteristic two this includes every Weierstrass
point and the purely two-primary non-Weierstrass examples of A09.
Constant extensions preserve that primary support. Other covers may
still work; no absence conclusion follows from this residual condition.

A30 refutes a proposed direct higher-genus extension of the
``every detecting quotient is good'' conclusion: a genus-three
rank-one curve has a detecting constant cyclic cubic quotient whose
relative class number is even. It does not exclude other good
quotients. The proposed unramified double-cover repair is unavailable
to the geometric prime-to-two fundamental group and does not preserve
the descending detector. A31 shows that the rank-one curve of A27
nevertheless has an exact signed-doubling theta witness of order seven.
All six of its original rational points permanently miss A29's
canonical detector, yet the geometric witness can be pointed after
a finite extension. Thus the sufficient tests have distinct limits;
their failure is not a nonexistence result.

The strongest new reduction of the characteristic-two genus-two
\emph{point-section problem} is A33: it suffices to prove the assertion
for ordinary genus-two curves with a Rosati-invariant integral
$\mathcal O_{\mathbf Q(\sqrt5)}$ action defined over the finite base field.
A cyclic quintic cover either has geometric new rank zero or gives such a
quotient; two reflection maps act on one common lifted collision.
This does not prove the separate literal primitive-cover problem.
A34 offers a further sufficient cover condition: for a constant pointed
odd cyclic cover of degree $n$, its actual rational two-primary Prym
exponent $M$ may be nontrivial if $n>(2M-1)^2$. An explicit degree-nineteen
cover after a degree-seven constant extension has $M=2$ and even relative
class number. The uniform Weil replacement for $M$ gives no improvement
on the older size test; the actual exponent or another geometric input
still has to be controlled.

A35 makes the purely two-primary canonical residual set a finite set:
all of its points lie over the field fixing geometric $J[4]$, whose degree
is at most six. For the ordinary real multiplication targets of A33 the
degree divides six. This does not treat nonzero odd-primary components
supported at primes where two has odd multiplicative order. A32 also
exhibits an ordinary twelve-dimensional new factor with odd rational
order over $\mathbf F_{256}$, losing that property exactly at extension
degrees divisible by seven. Geometric rank zero, rational unit order,
and bounded rational exponent therefore remain distinct sufficient tests.
The higher-genus and other-characteristic requirements of full A1 have
not been reduced to this genus-two real multiplication locus.


The real multiplication restriction is now verified to be an actual
proper class: the explicit ordinary curve (DRP.1) has geometric
endomorphism field with Rosati-fixed field $\mathbf Q(\sqrt2)$,
whereas the A13 curve has an integral Rosati-$\mathcal O_5$ action.
A37 closes the fixed-finite-support proposal inside this precise
smaller class: for every finite set of odd primes, including five,
there are densely many ordinary real multiplication curves with no
nonzero geometric theta point supported on those primes, with all
powers and products allowed. This does not exclude a prime depending
on the curve. A36 separately refutes the forced-five canonical
shortcut on an actual reflection target; its rational Jacobian order
is $2^{10}3^4$. The exact inhomogeneous constants have no five-primary
part there. A39 below supplies a three-primary detector on that particular target.
A universal construction for prescribed transferred points is still required.

A38 further rules out the unrestricted claim that ordinary
integral real multiplication forces a useful canonical prime at every
non-Weierstrass point: all $126$ such points of the A13 curve over
$\mathbf F_{128}$ have canonical class of exact order $4159$, and
$\operatorname{ord}_{4159}(2)=2079$ is odd. This curve still has the
geometric theta witness of A13. A40 now proves actual image membership for each prescribed geometric
point after a finite constant extension, using an inverse dihedral
construction for every ordinary real multiplication curve. Thus broad
canonical forcing fails even on actual images. A41 strengthens this to simultaneous failure at every lift and all
five reflection images of the same quintic cover, with at least eleven
split source points over $\mathbf F_{128}$. A43 proves that these
sources have two-rank one and that both of their Weierstrass points are
rational, split, and have only such bad canonical reflection images.
Thus these two qualifications do not repair the fixed-cover argument.
Whether a suitable choice of a different cover always succeeds remains open.

A39 gives a positive exact test on one earlier example: both
Weierstrass source points of A27 over $\mathbf F_{256}$ have useful
three-primary canonical detectors on the A36 quotient. A complete
$160\times155$ Riemann--Roch matrix is independently verified by
binary rank $1240$, and an explicit function identifies the image
where the canonical class is zero. Composite etale degrees $15$ or
$45$ suffice there; their geometric normal closures have odd order.
This example does not provide a universal supply theorem.


A42 makes the primary-projection alternative quantitative using one
fixed Lang cover. On an infinite collection of constant extensions, a
prescribed component has asymptotic fraction $1/\#G_S$ with an explicit
square-root error. In particular, after rationalizing a nonzero useful
odd-primary class, canonically detectable points are Zariski dense on
the Raynaud theta curve. The counting cover may have characteristic-primary
degree; it is not itself an allowable prime-to-characteristic
neighbourhood. Each resulting canonical detector is supplied separately
by A29. Density still does not put a good point into a prescribed finite
reflection fibre, and uncontrolled primary components prevent deduction
of whole-order primitive-orbit coverage. A41 and A43 show that this
distinction occurs on actual quintic transfers.


The special automorphism used in A43 is not a universal replacement
for cover selection. Its supplement shows that such an order-three
automorphism forces a nonzero order-three point on the actual Raynaud
theta divisor. A37 supplies ordinary real multiplication targets
without any such point, even geometrically. Also, varying A40's inverse
parameter changes the source curve, and its ramified degree-two
reflection map does not give a reverse admissible lifting operation.

A44 tests a controlled constant-field tower. For an abelian surface over
a field of characteristic two, every eligible odd prime occurring in
the tower of degrees $2^m$ divides the explicit nonzero integer (CT.2).
All points supported on those primes lie in the single finite field
of (CT.5), including arbitrary primary powers. At the same time,
the Lang lower bound and exact eventual two-primary growth prove that
there are infinitely many odd theta points in this tower. Thus
unbounded abundance alone does not supply unbounded eligible orders.
The finite candidate set need not be empty: (CT.15)--(CT.16) give
four exact order-three witnesses on A13. Universal success within that finite set is not settled by A44. Its
odd auxiliary-prime supplement proves that every compatible new order
has at most one new prime and a cofactor from an explicit finite list.
Poonen's bounded-index subgroup may contain multiple new primary
components, so that theorem does not supply the required prime-power
incidences. The supplement leaves their existence open.

A45 permits nonconstant geometric deck groups in the actual-exponent
criterion over the original pointed field. It proves separation whenever
$2\le g\le p$ and $n>(2M-1)^2$; in larger genus a collision forces a
separable elliptic map of degree $p^a\le g-1$, totally ramified at both
points. Primitive factors preserve their exponents at their own minimal
deck fields, but proper factors may grow at the full deck field. The
exact degree-$39$ cover has exponents $2$ and $16$, showing strictness
for this specified cover. No universal bound or supply follows.

A47 further allows a positive actual primitive exponent $M_n$ with
$n>(2^{\omega(n)}M_n-1)^2$, in the range $2\le g\le p$ or after
excluding its precise wild elliptic obstruction. For composite degree
this criterion can ignore large proper-conductor torsion, but the factor
$2^{\omega(n)}$ means that it is not uniformly stronger than A45.
A universal supply theorem for either sufficient inequality is still
missing. The literal rank-zero and rational-unit cover questions are
not settled by allowing a bounded positive exponent.

A48 supplies good unpointed twists of every genus-two cyclic cover
with constant deck group. This does not supply the specified pointing.
Exactly one twist is split above the prescribed point; every other twist
becomes isomorphic to that pointed model over an extension in which the
fibre is rational. Any original rational primitive characteristic-primary
torsion therefore persists. The affine-Frobenius calculation proves the
same obstruction for nonconstant deck groups. Varying the geometric cover
remains possible. The A47 supplement also proves that even a rounded
coarse Weil exponent bound cannot certify its numerical criterion.

A48's pullback supplement also shows why a nonordinary-cover iteration
does not repair an already failing cyclic cover: under a geometrically
connected prime-to-$p$ base change retaining its conductor, the old actual
primitive $p$-primary group injects into the new one over every common
finite field. Its geometric rank and rational exponent cannot decrease.
Raynaud's simultaneous pulled-back theta incidences do not imply
nilpotence in the larger coherent character spaces. Fresh cyclic covers
on the new base remain outside this obstruction.

The geometric A48 corollary realizes this distinction on actual curves:
after an odd Raynaud cover of degree $d=5\cdot3^a$, $13\le a\le65$,
all twelve primitive arrows of a cyclic degree-thirteen cover are singular,
while its geometric primitive two-rank lies between $12$ and $12(d-1)$.
This strengthens the evidence against the inference already refuted in A04;
it is not a new primitive-unit witness.

A49 makes the finite-quotient route exact. For an odd quotient killing
$[K_C-2P]$ whose exponent has unit group generated by $p$ and $-1$,
the point-image augmentation element is nonuniform exactly when some
cyclic quotient has a good actual primitive factor over the original
pointed field. The full odd quotient may succeed when a prime-level
size test fails. The two retained examples already have other witnesses,
so they do not establish a strict reduction of the full previous
universal existence gate. The quotient and nonuniformity must still
be supplied for each prescribed curve and point.

The variable-quotient supplement does supply nonuniformity when the
conductor is unrestricted. After fixing any finite excluded prime support
containing the canonical support, suitable constant extensions keep the
projection kernel fixed. The genus-two difference bound then supplies a
nonzero singleton coset and an empty coset. This gives one unit character
and its signed characteristic-power orbit, with pairwise coprime
conductors available. It does not supply the signed-unit condition or
control the other faithful character orbits of that conductor.

A50 checks whether a partial good packet can be used directly. A
cyclotomic norm bound forces integral projector coefficients to grow with
the required characteristic-primary precision. An exact abstract
$(\mathbf Z/8)[C_{31}]$ module has a missing ten-character signed packet
but minimum nonzero integral annihilator coefficient norm eight, too
large for the displayed degree comparison. This is not a geometric
counterexample. Neither this module calculation nor the inspected
nonabelian tower literature supplies the remaining actual cover bound.

A51 computes the exact norm-compatible point-difference orders under
prime-to-charac\-teristic refinement. The old torsion splits off, while a
new norm-kernel component remains. Norms to every proper cyclic quotient
kill the faithful product-conductor component, so the known packet
constraints do not control it by those maps. In a specified rank-one
tower, integral semilinear descent bounds the number of rational
characteristic-primary generators uniformly, but supplies no exponent
bound. Both arbitrary-exponent examples are explicitly abstract.

The full-fundamental-group point-section argument in
\cite[\S3, Theorem~3, pp.~134--135]{Faltings840} uses a level annihilating
the whole finite Jacobian group. Restricting its construction to levels
prime to the characteristic leaves precisely the characteristic-primary
residual. The inspected proof therefore does not repair this gate.
The alternative A05 proof retains its separate cuspidal/configuration
input.

The A51 supplements check two further repairs. A fixed nonzero
rational determinant in the character parameter would bound the actual
exponents, if the necessary geometric identification were proved.
Completed coefficients and finite matrix size alone do not do so.
The inspected skew-abelian tower result uses the tower prime, a growing
return product, and an integral preliminary statement whose natural-image
version fails for multiplication on an elliptic curve. Its separate
abstract results are not thereby refuted. In genus two and characteristic
two, full deck-Frobenius transitivity actually makes rational primitive
two-torsion nonzero whenever the primitive factor is ordinary. Its return
valuation remains the unknown quantity.

The coefficient hypotheses in the inspected Fourier--Mellin family
construction \cite[Theorems~1.3,1.5 and Definition~4.4]{EsnaultKerz2021}
and the unit-$L$-function result
\cite[Corollary~1.8]{EmertonKisin2001} do not supply that uniform
cross-prime determinant. In the latter result, the unit correction concerns
the $L$-function variable. It is not a rationality statement in the
varying finite-order character parameter.

For the iteration attempt, A51 proves that all geometric pointed
quintic/reflection choices have representatives over the constant extension
with degree support $\{2,3,5,13\}$, and gives an explicit common field at
each finite depth. That support does not obstruct good canonical points:
on every fixed target, a fixed Lang-cover count supplies densely many
points with a nonzero three-primary canonical component inside the same
extension. Incidence with the prescribed iterated orbit remains unproved.

The weighted-Weil supplement constructs actual ordinary principally
polarized abelian varieties with the transitive polynomial shape and one
rational characteristic-primary generator, but exponentially large exponent.
It does not realize the required cyclic deck action or Jacobian cover.
A52 permits any integral annihilator with nonzero faithful evaluation and
small enough coefficient norm. Its fixed-field lattice bound depends on
the actual ideal index $W_r$, which remains unbounded by the proved inputs.
In the fully transitive case it recovers exactly the earlier exponent test.
The actual orbit-module norm lower bound and integral trace descent exclude
an improvement from rational generator counts or inflation of an older
principal relation alone. No actual witness outside the full preceding
union of criteria, and no universal $W_r$ bound, is proved.

The final actual-point calculation gives a nonzero regular logarithmic
differential of trace zero when the compatible point-class order increases.
The two forced wild fibres cancel the degree term in the elementary
Riemann--Hurwitz lower estimate, so neither calculation bounds the order.
Finitely many good signed packets can also be preserved during further
variable-conductor supply, by a common arithmetic progression of constant
degrees. On an actual common cover the norm operators still act block by
block. Even all line quotients of $(\mathbf Z/\ell)^2$ leave a missing
signed orbit inside one rational cyclotomic block uncontrolled by the other
line norms. No universal bound against coefficient cancellation, and no
new small operator on the actual point class, is claimed.

The approaches actually tested were: fixed-base Hooley density;
Goldfeld's average over bases; Lenstra's fixed finite Chebotarev
constraints and positivity; prime-to-characteristic torsion orders
from Pop--Sa\"idi; a single singular Frobenius arrow in higher genus;
scalar-Frobenius packet counts; generic ordinary and fixed-rank
Prym theorems with finite-type specialization; a direct degree-$13$
theta certificate; prime-power character blocks; primary torsion
projections on theta; the Lang quotient and Weil bound; growth of
fixed-prime primary torsion; the composite primitive projector;
an exact order-$15$ search and certificate; the Hooley--Lenstra
family $n=3q$; the elliptic-exception rank calculation; finite
prime-support torsion via Frobenius translation orbits; the
all-three-primary obstruction; constant twists with an unspecified
split base point; the parity augmentation test; its all-extension rank-one obstruction;
inverse-character canonical detection and its prime-power refinement;
two-reflection transfer to the ordinary real multiplication locus;
actual rational Prym exponent with a degree-nineteen certificate;
the complete two-primary theta locus through geometric four-torsion;
fixed Lang covers with prescribed primary components; the exact
rank-one and two-Weierstrass analysis of the inverse family;
the attempted extension of its order-three automorphism;
finite eligible-prime support in a characteristic-power tower;
all-conductor Fourier inversion on an odd rational quotient;
nonconstant-deck geometric towers and the scope of class-number stability;
weights and principal polarization; fixed-field annihilator lattices;
trace relations and exact orbit-module norm bounds;
and the independent cuspidal route.
Their respective endpoints are recorded in A02--A52 and the
retained B.28--B.29 pointing audit. Poonen's
projection theorem supplies an order $\ell^r m$, with uncontrolled
prime factors of $m$; multiplying by $m$ moves theta, and passing to
an intermediate cover leaves proper-conductor character blocks.
Those are failures of those two strategies. A15 bypasses the
proper-conductor factors by a product projector on its stated
geometric locus, and A18 removes its geometric exception in
characteristic two. What remains is the same-subgroup orbit-coverage
condition, or the stronger single-point unit-group condition.
A17 gives density $\mathcal A/2$ to the candidate primes
$q\equiv5\pmod8$ for $n=3q$, conditionally on M, without proving
theta incidence for any particular curve. A12 gives an unconditional
curve-dependent primary-torsion test; A14 proves that enlarging the
constant field at one fixed prime cannot eventually make that
particular sufficient test succeed. Neither statement is a theorem
of nonexistence of theta torsion. The recovered Successor A.10
rational prime-cover criterion is separately distinguished in A14.
A19 now makes all primary exponents a finite exact search for
each prescribed finite odd prime support. A20 proves that even
all powers of three fail on the rank-one curve of A09, which
nevertheless has its order-thirteen witness. The prime supports
remain unbounded; no theorem proves that one of these finite
searches succeeds for every positive-rank curve.
The B.28--B.29 rational-twist alternative supplies a good split
point without prescribing it. Its good-point density estimate
does not force such a point into an arbitrarily prescribed fibre,
and extending a good twist can destroy its rational torsion
condition. These failures are recorded separately from geometric
primitive-rank-zero incidence.
The cuspidal route
settles point separation, while the scalar route refutes the
universal auxiliary existence assertion. None of the other routes
silently fills the curve-dependent incidence gap.

\paragraph{U02: the literal primitive-cover problem.}
ANABELIAN\_A1 B.24 allows the degree-one cover with convention
$P_1=J_C$. The B.23 counterexample has $p$-rank zero and is therefore
covered by that convention. The problem for positive $p$-rank
remains open in the recovered run. Similarly ANABELIAN\_A1 B.5 and
B.32 ask for prescribed-point rational class-number witnesses. A29
now settles the stated nonzero canonical-primary component cases,
while their universal existence remains open.
A bound for a prime in a number-field Frobenius class does not
produce the required covering geometry or local unit conditions.
No assertion is made that these separate questions follow from A05.

\paragraph{U03: complete source coverage.}
The two missing registers account for 390 full statements. Only
47 distinct A1-successor IDs and 50 distinct GT-tower IDs are visible
in the supplied summaries and flag excerpts. The missing LM7
proof files account for another 209 proof blocks. Exact-title,
variant-title and content searches recovered four relevant archives,
but not the two missing sources. A final prior-context recovery
attempt was unavailable because that retrieval service was disabled.
The next concrete input is the complete
\texttt{ANABELIAN\_A1\_SUCCESSOR} register/paper and the complete
\texttt{GT\_TOWER\_CHECK} register/paper, together with the omitted
LM7 proof sections. Their absence is an evidentiary gate, not a
mathematical theorem beyond GRH.

\paragraph{U04: infinite towers and unspecified finite data.}
For a recovered row with finite affine monodromy, G01 gives a
prime bound after the actual finite image, cocycle, precision,
and excluded primes are fixed. It does not bound the precision
needed to separate an infinite object, the height of a Hilbert
specialization, or a Baire-category construction. A finite test of
full adic semisimplifications is available under the hypotheses of
G03, with exact computable traces as an additional algorithmic input.
Its applicability to an unspecified row must be checked separately.

\paragraph{U05: scope of the effective open-image bounds.}
The former descent gate for LM8 A.32--A.33 is closed by their actual
models $F_{\pm6}/\mathbb Q$ and the checked refinement G06.
The conditional Mayle--Wang bound applies over $\mathbb Q$; after
restriction to the proof's number field $L$, G06 proves an index
bound $[L:\mathbb Q]$, not necessarily full image.
G-LOMBARDO is an unconditional explicit adelic index bound over
arbitrary number fields. The hypotheses must still be checked for
any other factor; CM factors have CM monodromy rather than full
$\operatorname{GL}_2$ monodromy. No corresponding general abelian-variety
bound has been asserted.

\paragraph{U06: boundaries of verification.}
The complete new proofs were read, and every included computational
certificate reports the exact execution actually performed. No
formal proof-assistant verification of this paper, reconstruction
of all earlier profinite elements, replay of every historical
computer search, or verification of Masullo's entire draft is
claimed. The check table distinguishes whole-proof reading,
independent arithmetic reimplementation, retained-verifier replay,
source-statement checks and carried background. A same-model second
or third reader is not presented as verification by a different
model or by a theorem prover.
